% FrontBackMatter/Abstract.tex
\addchaptertocentry{Хураангуй}

\vspace*{-1cm}
\noindent\rule{\textwidth}{0.4pt}
\vspace{0.8cm}

\begin{center}
{\Large\textbf{\ttitle}\par}
\bigskip
{\scshape\Large \deptname\par}
\bigskip

{\large \shortname\par}
{\Large\textbf{Хураангуй}\par}

{\normalsize \par}
\end{center}
\vspace{1cm}


\bigskip


Монгол бичиг нь Монгол үндэстний соёл, түүхийн үнэт өвийг хадгалсан уламжлалт бичиг юм. Гэвч орчин үед Монгол бичгийг сурах, хэрэглэх боломж хязгаарлагдмал байгаа нь залуу үеийнхэн уламжлалт бичгээ мартаж, соёлын өв алдагдах эрсдэлийг бий болгож байна. Мөн одоо байгаа Монгол бичиг сургах материалууд нь ихэвчлэн уламжлалт хэлбэртэй байдаг тул залуу үеийнхний анхаарлыг татахгүй, интерактив бус байна.

Энэхүү судалгааны ажлын зорилго нь Монгол бичгийг сонирхолтой, хялбар, орчин үеийн технологийн тусламжтайгаар суралцах боломжийг олгох \\textbf{Egüü} гар утасны аппликейшныг хөгжүүлэхэд оршино. Систем нь React Native дээр суурилсан cross-platform апп бөгөөд Firebase/Firestore backend, TensorFlow OCR (Монгол бичиг таних), flashcard сургалтын систем, олон хэлний дэмжлэг (i18n), dark/light theme, streak tracking, daily word зэрэг орчин үеийн функцуудыг агуулна.

Төслийн үр дүнд хэрэглэгчид Монгол бичгийн үсгүүдийг зураг, фонетик, орчуулга, тодорхойлолтын хамт суралцах, өөрийн гараар бичсэн үгийг таниулах, өдөр бүрийн шинэ үг авах, дуртай үгсээ хадгалах зэрэг боломжуудыг олж авна. Энэ нь Монгол бичгийг сурах процессыг сонирхолтой, үр дүнтэй болгож, уламжлалт соёлын өвийг хадгалахад хувь нэмэр оруулах юм.



\vfill